%% kapa: exam=1 round=35 year=2024 period=1 form=A subject=economics subject_ko=경제학원론 seq=21 number=61
%% source: Q-Net 감정평가사 기출문제 (cst00309) artlSeq=5213229 file=raw/1008/5213229/1교시 문제지 원본.pdf
%% conversion: kordoc 4.24.51 -> tools/gen1.py (auto); review=proofread-sonnet; flags=none
\begin{question}{61}
두 재화 $X$, $Y$만을 생산하는 A국의 2022년과 2023년의 생산량과 가격이 아래와 같다. 2023년의 전년대비 (ㄱ)\uline{경제성장률(실질GDP증가율)}과 평균적인 가계의 소비조합이 $X$재 2단위, $Y$재 1단위일 때 (ㄴ)\uline{소비자물가상승률}은? (단, 기준연도는 2022년이다.)
\begin{center}\small\begin{tabular}{|c|c|c|c|c|}\hline
\multirow{2}{*}{연도} & \multicolumn{2}{c|}{$X$} & \multicolumn{2}{c|}{$Y$} \\ \cline{2-5}
 & 수량 & 가격 & 수량 & 가격 \\ \hline
2022년 & 100 & 10 & 80 & 50 \\ \hline
2023년 & 100 & 15 & 100 & 40 \\ \hline
\end{tabular}
\end{center}
\begin{choices}
  \item ㄱ: 10 \%, ㄴ: 0 \%
  \item ㄱ: 10 \%, ㄴ: 10 \%
  \item ㄱ: 20 \%, ㄴ: -10 \%
  \item ㄱ: 20 \%, ㄴ: 0 \%
  \item ㄱ: 25 \%, ㄴ: 10 \%
\end{choices}
\end{question}
\answer{4}
